
Spurious correlations degrade worst-group accuracy while maintaining high average accuracy. Existing detection methods require knowing which spurious feature to search for. Robust learning approaches demand group annotations (GroupDRO, JTT) or operate post-hoc on embeddings (SCER). Gradient attribution methods fail for unknown spurious correlations.

We proposed a gradient-based framework that detects and mitigates spurious correlations via gradient abnormality, extending GAIA from distribution shift to subpopulation shift. The methodology operates on gradient process (abnormality, not attribution) and provides unified detection-mitigation in a single pipeline.

\textbf{Validated contributions:}

Detection (h-e1): Gradient abnormality pipeline differentiates minority vs majority patterns in synthetic experiments (divergence 0.30, p<0.0001, Cohen's d=198.75). First application of GAIA to spurious correlation detection.

Mechanism (h-m-integrated): Causal validation via synthetic tests shows strong correlation between GAIA divergence and worst-group accuracy (|$\rho$|=0.975, p<0.001), background augmentation reduces GAIA-Z by 39.4\% (p<0.001, d=2.96), supporting spurious-conflict hypothesis.

Mitigation (h-m-mitigate): Spatial gradient regularization improves worst-group accuracy by +23pp on MNIST+Color toy dataset (78\% vs 55\%) in proof-of-concept experiments without catastrophic accuracy drop.

\textbf{Limitations:} All validation from synthetic data (h-e1, h-m-integrated) or minimal proof-of-concept (h-m-mitigate). Real Waterbirds experiments require GPU compatibility resolution (PyTorch 2.1+ for H100 sm\_90 support) or CPU training (~70-90 GPU hours). Empirical claims (real minority gradients exhibit abnormality, real Waterbirds worst-group accuracy improvement) pending full-scale validation.

\textbf{Contribution tier:} Tier 3 (methodology validated) with clear path to Tier 2 (real empirical validation). Synthetic validation proves code correctness and mechanism plausibility; real validation required for empirical claim verification.

\textbf{Immediate future work:} Real Waterbirds detection/mitigation experiments, Full 5-seed MNIST validation, GroupDRO baseline comparison. \textbf{Extensions:} Multi-dataset generalization (CelebA, UrbanCars), joint gradient-embedding regularization, automatic hyperparameter selection.

The gradient space encodes training dynamics --- what models rely on during learning. Our work demonstrates gradient abnormality as a viable paradigm for spurious correlation detection, complementing embedding-based approaches and bypassing attribution ineffectiveness. Methodology validated; empirical validation in progress.
